\section{Compile-Time Instrumentation}

As already detailed (Section \ref{sec:compandinst}), visibility into the target application is achieved through compile-time instrumentation. In our solution, the instrumentation logic is integrated into the LLVM compiler infrastructure, specifically extending the existing LLVM mode's PCGUARD pass. 

The instrumentation process is explicitly designed for indirect branches and consists of three primary phases: target identification, metadata allocation, and callback injection.

\paragraph{Target Identification} The first phase of the instrumentation pass involves traversing the LLVM Intermediate Representation to identify instructions capable of dynamic control flow transfers. 

During the traversal of each basic block, the pass filters the instruction stream to isolate two specific LLVM IR instructions: 
\begin{itemize}
    \item \texttt{IndirectBrInst}: Represents an indirect branch instruction --- e.g., computed \texttt{goto} statements and jump tables.

    \item \texttt{CallInst}: Represents a function call. In this case, the pass further verifies that the target operand is an indirect function pointer call or a virtual method dispatch.
\end{itemize}

To maintain feature parity with the broader AFL++ ecosystem, this identification phase respects the coverage boundaries set for standard edge coverage. Fuzzing campaigns can utilise allowlists and denylists to constrain instrumentation to specific source files or functions.

\paragraph{Guard Array Architecture} Once an indirect instruction has been verified and selected for instrumentation, the compiler allocates a slot to track its execution. To track each indirect branch source, our architecture utilises a metadata structure referred to as "Guard Array".

For every identified indirect instruction, the compiler pass allocates a 32-bit integer slot --- a \texttt{GuardPtr} --- into a dedicated contiguous ELF section named "\texttt{\_\_sancov\_indir\_guard}" (mimicking the operation of PCGUARD mode).

As the compiler processes each translation unit, it continuously appends 32-bit slots into this section. The resulting linked binary will contain a sequential array of uninitialized metadata slots, corresponding one-to-one to every instrumented indirect branch source. At runtime, these slots will be populated with unique sequential identifiers.

\paragraph{Callback Injection} The final phase of static instrumentation involves modifying the target's executable code to report the control flow transition. This requires injecting a C function callback into the LLVM IR.

Modifying the IR of a basic block while actively iterating over its instructions risks invalidating iterators and triggering undefined behaviour within subsequent analysis passes. To circumvent this, a two-phase approach is taken: 
\begin{enumerate}
    \item \textit{Collection Phase:} The pass first traverses the CFG, applying the identification and allowlist logic to build a set of valid targets.

    \item \textit{Mutation Phase:} The pass performs the necessary IR modification on the collected set of targets.
\end{enumerate}

For each target instruction in the mutation set, the compiler injects a call to \texttt{\_\_afl\_trace\_indir} immediately preceding the indirect branch. The callback takes two arguments: the address of the corresponding \texttt{GuardPtr} slot and the target address of the indirect branch, which the program has just evaluated.